% =====================================================================
%  sections/99-tail.tex — bibliografia, spisy i dodatki.
%  Kolejność bloków poniżej odpowiada wymogom wydziału; redagujesz tu
%  głównie Dodatek A (opis tego, co trafia na płytę/archiwum).
% =====================================================================

% Bibliografia: \printbibliography składa wpisy z bibl.bib (biblatex).
% hbadness/hfuzz wyciszają ostrzeżenia o składu — wpisy z DOI i URL-ami
% łamią zasady justowania nieodwołalnie i zaśmiecałyby log.
{
\hbadness=99999
\hfuzz=\maxdimen
\addcontentsline{toc}{section}{Bibliografia}
\printbibliography[title = {Bibliografia}]
}

% Spis rysunków generuje się sam z \caption (styl: tocloft w preambule).
\newpage
\listoffigures

\newpage
\phantomsection % Ważne dla poprawnych linków w PDF (pakiet hyperref)
\addcontentsline{toc}{section}{Spis algorytmów} % Ręczne dodanie wpisu
\listofalgorithms

\pagebreak
\phantomsection % Ważne dla poprawnych linków w PDF (pakiet hyperref)
\addcontentsline{toc}{section}{Dodatek A. Archiwum plików}
\pagestyle{empty}
\section*{Dodatek A. Archiwum plików}
\subsection*{Zawartość archiwum}
W archiwum umieszczono kod źródłowy tego dokumentu (pracy inżynierskiej),
kod źródłowy biblioteki,
przykładów jej wykorzystania oraz pliki wykonywalne na system operacyjny Linux.
Dołączone zostały również karta dyplomowa oraz oświadczenie dot. GenAI.
\subsection*{Uruchamianie przykładów}
Zaleca się uruchamianie programów z poziomu terminala. Wymagany jest procesor
z obsługą AVX512. Należy nadać plikowi uprawnienia do uruchamiania (jeśli ich nie ma):
\begin{verbatim}
chmod +x program
\end{verbatim}
Należy następnie go uruchomić:
\begin{verbatim}
./program
\end{verbatim}
W przypadku chęci samodzielnej kompilacji i uruchomienia przykładów, należy
przejść do katalogu \texttt{mlrs} i użyć menedżera pakietów
\texttt{cargo}:
\begin{verbatim}
cargo run --release --bin [nazwa przykładu]
\end{verbatim}
kompilacja jest gwarantowana w przypadku użycia kompilatora w wersji
\begin{verbatim}
rustc 1.92.0-nightly (b6f0945e4 2025-10-08)
\end{verbatim}
UWAGA: kompilacja jest gwarantowana tylko na systemie linux debian 13 (trixie).
Biblioteka jak i przykłady powinny kompilować się i działać na innych systemach operacyjnych
-- autor nie daje jednak takiej gwarancji.

\newpage
%\pagestyle{plain}
%\pagenumbering{Roman}
%\setcounter{page}{1}
\section*{Dodatek B: Konwencje notacyjne pseudokodu}
\addcontentsline{toc}{section}{Dodatek B: Konwencje notacyjne pseudokodu}

W niniejszej pracy stosuje się notację pseudokodu łączącą elementy programowania funkcyjnego oraz imperatywnego, odpowiadając hybrydowemu charakterowi biblioteki napisanej w języku Rust.

\section*{B.1 Notacja ogólna}

\subsection*{Przypisania i mutacje}
\begin{itemize}
	\item $x \leftarrow \text{wartość}$ -- przypisanie wartości do zmiennej
	\item $L.\texttt{field}$ -- dostęp do pola struktury/obiektu
	\item $\text{list}[i]$ -- indeksowanie tablicy/wektora (numeracja od 0)
	\item $\text{list}.\text{append}(x)$ -- dodanie elementu na koniec listy (mutacja)
\end{itemize}

\subsection*{Struktury kontrolne}
\begin{itemize}
	\item \texttt{ForEach} $x \in X$ -- iteracja po elementach kolekcji
	\item \texttt{While} warunek -- pętla z warunkiem
	\item \texttt{If}/\texttt{eIf}/\texttt{Else} -- instrukcje warunkowe
	\item \texttt{Return} wartość -- zwrócenie wyniku funkcji
\end{itemize}

\section*{B.2 Notacja funkcyjna}

\subsection*{Operacje wyższego rzędu}
\begin{itemize}
	\item $\text{map}(f, xs)$ -- aplikacja funkcji $f$ do każdego elementu listy $xs$
	\item $\text{zipWith}(f, xs, ys)$ -- sparowanie elementów dwóch list i aplikacja funkcji binarnej:
	      $$\text{zipWith}(f, [x_1, x_2, \ldots], [y_1, y_2, \ldots]) = [f(x_1,y_1), f(x_2,y_2), \ldots]$$
	\item $\text{zip}(xs, ys)$ -- sparowanie elementów: $\text{zip}([x_1, x_2], [y_1, y_2]) = [(x_1,y_1), (x_2,y_2)]$
	\item $\text{fold}(f, acc, xs)$ -- redukcja listy z akumulatorem początkowym
\end{itemize}

\subsection*{Równoległość}
\begin{itemize}
	\item $\text{par\_zipWith}(f, xs, ys)$ -- równoległa wersja \texttt{zipWith}
	\item $\text{parallel\_map}(f, xs)$ -- równoległa aplikacja funkcji
	\item Indeks dolny $_{\parallel}$ lub adnotacja \texttt{(równolegle)} -- oznaczenie równoległego wykonania
\end{itemize}

\subsection*{Notacja zbiorowa (list comprehension)}
\begin{itemize}
	\item $[\, f(x) \mid x \in xs, P(x) \,]$ -- lista elementów $f(x)$ dla $x$ spełniających predykat $P$
	\item $\{\, a + b \mid (a,b) \in \text{zip}(xs, ys) \,\}$ -- zbiór sum sparowanych elementów
\end{itemize}

\section*{B.3 Operacje algebraiczne}

\subsection*{Operacje macierzowe}
\begin{itemize}
	\item $A \cdot B$ lub $\text{matmul}(A, B)$ -- mnożenie macierzy
	\item $A^T$ lub $\text{transpose}(A)$ -- transpozycja macierzy
	\item $A \odot B$ lub $\text{elem\_mul}(A, B)$ -- mnożenie element-wise (Hadamard)
	\item $A \oplus \mathbf{b}$ lub $\text{broadcast\_add}(A, \mathbf{b})$ -- dodawanie z rozszerzaniem (broadcasting)
\end{itemize}

\subsection*{Funkcje agregujące}
\begin{itemize}
	\item $\sum_{\text{axis}} A$ lub $\text{sum\_axis}(A, \text{axis})$ -- suma wzdłuż wymiaru
	\item $\text{sum\_rows}(A)$ -- suma wierszy (zwraca wektor kolumnowy)
\end{itemize}

\subsection*{Funkcje aktywacji}
\begin{itemize}
	\item $\sigma(z)$ -- funkcja aktywacji (zazwyczaj sigmoidalna)
	\item $\sigma'(z)$ lub $\text{sigma\_derivative}(z)$ -- pochodna funkcji aktywacji
	\item Aplikacja na macierz: $\sigma(Z)$ oznacza element-wise: $[\sigma(Z)]_{ij} = \sigma(z_{ij})$
\end{itemize}

\section*{B.4 Konwencje specyficzne dla tensorów}

\begin{itemize}
	\item $T.\text{shape}$ -- wektor wymiarów tensora, np. $[C, W, H, B]$
	\item $T.\text{strides}$ -- wektor kroków (strides) w pamięci
	\item $T.\text{data}$ -- surowy wektor danych tensora
	\item $T.\text{iter}()$ -- iterator po elementach w kolejności logicznej
	\item $T.\text{iter\_indices}()$ -- iterator po wielowymiarowych indeksach
	\item $T[idx]$ -- dostęp do elementu przez indeks wielowymiarowy
\end{itemize}

\section*{B.5 Notacje pomocnicze}

\begin{itemize}
	\item $|xs|$ -- długość listy/wektora
	\item $xs \times ys$ -- iloczyn kartezjański (wszystkie pary)
	\item $\text{reversed}(xs)$ -- odwrócenie kolejności elementów
	\item \texttt{// komentarz} lub \tcp{komentarz} -- komentarz wyjaśniający
\end{itemize}

\section*{B.6 Przykłady równoważnych zapisów}

Dla ilustracji, następujące zapisy są równoważne:

\begin{align*}
	\text{zipWith}(+, xs, ys) & \equiv [\, x + y \mid (x,y) \in \text{zip}(xs, ys) \,]   \\
	                          & \equiv \text{map}((x,y) \mapsto x+y, \text{zip}(xs, ys)) \\
	A \cdot B                 & \equiv \text{matmul}(A, B)                               \\
	\sigma(Z)                 & \equiv \text{map}(\sigma, Z) \text{ (element-wise)}
\end{align*}

\section*{B.7 Uwagi implementacyjne}

\begin{itemize}
	\item Operacje oznaczone jako \emph{równoległe} odpowiadają użyciu biblioteki Rayon w implementacji Rust
	\item Notacja funkcyjna (zipWith, map) odpowiada stylowi \emph{iterator chaining} w Rust
	\item Operacje macierzowe są delegowane do biblioteki \texttt{faer}
	\item Wszystkie algorytmy zakładają układ pamięci \emph{column-major} (zerowy wymiar ciągły)
\end{itemize}

\section*{B.8 Obsługa błędów}

W języku Rust funkcje zwracają typ \texttt{Result<T, E>}, będący monadą reprezentującą albo sukces (\texttt{Ok(value)}), albo błąd (\texttt{Err(error)}). W pseudokodzie:

\begin{itemize}
	\item $\text{Ok}(x)$ -- operacja zakończona sukcesem, zwraca wartość $x$
	\item $\text{Err}(e)$ -- operacja nie powiodła się, zwraca błąd $e$
	\item \textbf{panic:} -- nieodwracalny błąd (odpowiednik \texttt{panic!} w Rust)
	\item $\bot$ -- brak wartości (odpowiednik \texttt{None} w \texttt{Option<T>})
\end{itemize}

Algorytmy wykorzystujące walidację wejścia (np. sprawdzanie granic) zwracają \texttt{Result}. Algorytmy, w których błąd jest niemożliwy lub nieistotny, pomijają tę notację dla czytelności.
